\chapter{Bowel Obstruction}

\section*{Definition and Overview}
Bowel obstruction is a blockage that prevents food, fluid, and gas from moving through the intestine. It can affect the small bowel or the large bowel, can be complete or partial, and can develop suddenly (acute) or gradually (chronic). A bowel obstruction is a medical emergency: if the bowel becomes completely blocked and its blood supply is compromised, the bowel wall can become ischaemic, die, or perforate, leaking bowel contents into the abdominal cavity and causing peritonitis and sepsis. Treatment usually requires hospital admission, and the outcome depends heavily on the cause and on how quickly treatment begins. With prompt care, many obstructions resolve without surgery, and even when surgery is needed the outlook is generally good.

\section*{Epidemiology}
Bowel obstruction is one of the most common surgical emergencies, accounting for a substantial number of hospital admissions worldwide. In adults, the leading cause of small bowel obstruction is adhesions -- bands of scar tissue from previous abdominal or pelvic surgery -- which carry a lifelong risk that rises with each operation. The most common cause of large bowel obstruction is cancer, usually of the colon. Obstruction can occur at any age: in children, the important causes are intussusception, hernias, and congenital abnormalities such as malrotation. Older adults are especially vulnerable to obstruction from tumours, faecal impaction, and volvulus, and people who have had abdominal surgery, radiotherapy, or inflammatory bowel disease remain at risk throughout their lives.

\section*{Aetiology and Pathophysiology}
The causes of bowel obstruction are many, and are best divided by mechanism:
\begin{itemize}
    \item \textbf{Adhesions:} bands of scar tissue from previous surgery trap, kink, or compress a loop of bowel.
    \item \textbf{Hernias:} a loop of bowel becomes trapped in a weakness in the abdominal wall and cannot slide back (incarceration), and its blood supply may then be cut off (strangulation).
    \item \textbf{Tumours:} bowel cancer or other growths narrow the lumen and eventually block it.
    \item \textbf{Strictures:} narrowed segments of bowel caused by Crohn's disease, diverticulitis, or previous radiotherapy.
    \item \textbf{Volvulus:} the bowel twists on its mesentery, obstructing it and simultaneously compromising its blood supply.
    \item \textbf{Intussusception:} part of the bowel telescopes into the segment ahead of it, most often in children under two.
    \item \textbf{Faecal impaction:} hard stool packing the bowel, usually in older, immobile, or constipated people.
    \item \textbf{Paralytic ileus:} the bowel stops contracting after surgery, illness, or certain medicines, without any physical blockage.
    \item \textbf{Mechanism:} whatever the cause, the bowel above the blockage fills with gas and fluid and distends, causing pain and vomiting. If the blood supply is cut off, the bowel wall becomes ischaemic, may perforate, and can leak bacteria and bowel contents into the peritoneal cavity, producing peritonitis and sepsis.
\end{itemize}

\section*{Clinical Features}
The classic symptoms of bowel obstruction come on over hours to days:
\begin{itemize}
    \item Abdominal pain in cramping waves, often centred around the navel
    \item Vomiting, initially of stomach contents and bile, and later faeculent in appearance
    \item Abdominal swelling and distension
    \item Constipation and the inability to pass wind, although this may be absent early in a partial obstruction
    \item Loud, gurgling bowel sounds in the early stages, giving way to silence as the bowel loses function
    \item Fever, a rapid heart rate, and severe constant pain may signal that the bowel is dying or has perforated
    \item In children with intussusception, episodic crying and drawing up of the knees, with stool likened to redcurrant jelly
\end{itemize}

\section*{Differential Diagnosis}
Other conditions can mimic a bowel obstruction:
\begin{itemize}
    \item \textbf{Severe constipation or faecal impaction:} without a true mechanical blockage.
    \item \textbf{Paralytic ileus:} the bowel stops moving after surgery, illness, or medication.
    \item \textbf{Appendicitis, diverticulitis, or other causes of peritonitis:} which can present with pain, vomiting, and abdominal tenderness.
    \item \textbf{Biliary or renal colic:} intense colicky pain from gallstones or kidney stones.
    \item \textbf{Gynaecological emergencies:} ovarian torsion or cysts in women of reproductive age.
    \item \textbf{Food poisoning or gastroenteritis:} severe cramping and vomiting without obstruction.
    \item \textbf{Irritable bowel syndrome:} severe bloating and altered bowel habit.
\end{itemize}

\section*{When to Seek Medical Care}
Bowel obstruction cannot be treated at home. Go to an emergency department immediately if you have the combination of cramping abdominal pain, vomiting, abdominal swelling, and the inability to pass wind or stool. Do not eat, drink, or take laxatives, and seek help without delay, because the danger increases the longer the bowel remains blocked.

\textbf{Call 000 (or local emergency services) for:}
\begin{itemize}
    \item Severe, constant abdominal pain that is getting worse
    \item Vomiting together with a swollen abdomen
    \item Signs of shock: faintness, a rapid heart rate, cold clammy skin, or confusion
    \item A tender, rigid abdomen that hurts to touch (possible perforation)
    \item High fever with abdominal pain
\end{itemize}

\section*{Diagnosis and Investigations}
\begin{itemize}
    \item \textbf{History and examination:} questions about previous surgery, hernias, cancer, and inflammatory bowel disease; examination for distension, tenderness, a palpable lump, and altered bowel sounds.
    \item \textbf{Blood tests:} full blood count, electrolytes, renal function, and inflammatory markers to assess dehydration, sepsis, and organ compromise.
    \item \textbf{Abdominal X-ray:} may show loops of bowel distended with gas and characteristic fluid levels, and helps identify perforation.
    \item \textbf{CT scan:} the most useful investigation, localising the site of blockage, identifying the cause, and showing whether the bowel wall is at risk.
    \item \textbf{Ultrasound:} valuable in children, particularly for confirming intussusception.
    \item \textbf{Water-soluble contrast studies:} can outline the site of a partial blockage and, in adhesive obstruction, sometimes help it to resolve.
    \item \textbf{Nasogastric tube:} passed to the stomach to decompress it, relieve vomiting, and allow monitoring of fluid loss.
\end{itemize}

\section*{Management}
Treatment depends on the cause, severity, and whether the blood supply is at risk:
\begin{itemize}
    \item \textbf{Admission and observation:} most people are admitted for monitoring, pain relief, and supportive care.
    \item \textbf{Bowel rest:} nothing by mouth, with a nasogastric tube to drain stomach contents and relieve pressure.
    \item \textbf{Intravenous fluids:} to correct dehydration and electrolyte disturbance from vomiting and fluid sequestration.
    \item \textbf{Conservative management:} many adhesive obstructions resolve with supportive care alone within a few days.
    \item \textbf{Surgery:} needed if the blockage does not settle, if there are signs of strangulation or perforation, or if the cause is a tumour; the obstructed segment may be resected and the bowel joined or brought out as a stoma.
    \item \textbf{Endoscopic treatment:} a self-expanding stent can sometimes be placed to open an obstructed bowel, particularly in advanced cancer or when surgery is not suitable.
    \item \textbf{Enemas and flatus tubes:} used for faecal impaction and for some cases of sigmoid volvulus, which may also be untwisted endoscopically.
    \item \textbf{Treating the cause:} hernia repair, resection of tumours, and management of strictures or inflammatory disease.
\end{itemize}

\section*{Complications}
The most feared complications arise when the blood supply to the bowel is compromised or the bowel perforates:
\begin{itemize}
    \item Strangulation, in which the blood supply is cut off and the bowel wall dies
    \item Perforation of the bowel, causing peritonitis and abscess formation
    \item Sepsis and septic shock from leakage of bowel contents into the abdomen
    \item Dehydration and acute kidney injury from vomiting and fluid shifts
    \item Short bowel syndrome if large amounts of bowel must be removed
    \item Recurrence of obstruction, particularly when adhesions are the cause
    \item Complications of surgery, including infection, bleeding, anastomotic leaks, and the formation of further adhesions
\end{itemize}

\section*{Prognosis and Outcomes}
The outlook depends on the cause and the speed of treatment. Many adhesive obstructions settle with conservative care within a few days, and surgery is usually successful when it is performed promptly. The main danger is delay: if the bowel is strangulated or perforated before treatment, the risk of serious complications and death rises sharply, and urgent surgery is associated with higher morbidity than elective procedures. Most people make a full recovery, but the outlook is more guarded when obstruction is due to advanced cancer, and some people experience recurrent episodes. After treatment, most people return to a normal diet and lifestyle, and those with an underlying cause such as a hernia or tumour are cured of the obstruction by addressing it.

\section*{Prevention and Self-Care}
\begin{itemize}
    \item Seek medical care early for any combination of severe cramps, vomiting, and abdominal swelling.
    \item After abdominal surgery, follow advice about early mobilisation and avoiding constipation, which reduce the risk of adhesive obstruction.
    \item Eat a high-fibre diet and drink enough fluids to avoid severe constipation and impaction.
    \item Manage conditions such as Crohn's disease and diverticulitis under medical supervision.
    \item Report a new bulge in the groin or abdomen that becomes painful or cannot be pushed back, since it may be an incarcerated hernia.
    \item Take laxatives only as directed, since overuse can contribute to bowel problems.
    \item If you are at risk of recurrence, know the early warning symptoms and do not wait to seek help.
\end{itemize}

\section*{Sources and Further Reading}
The following sources provide reliable, current information on bowel obstruction and its management.
\begin{itemize}
    \item NHS. \textit{Bowel obstruction.} https://www.nhs.uk/conditions/bowel-obstruction/
    \item Mayo Clinic. \textit{Intestinal obstruction.} https://www.mayoclinic.org/
    \item Centers for Disease Control and Prevention. \textit{Digestive diseases.} https://www.cdc.gov/
    \item World Health Organization. \textit{Digestive health.} https://www.who.int/
    \item MSD Manual. \textit{Intestinal obstruction.} https://www.msdmanuals.com/
    \item MedlinePlus. \textit{Bowel obstruction.} https://medlineplus.gov/
\end{itemize}
